\documentclass[11pt]{article}

\usepackage[T1]{fontenc}
\usepackage{lmodern}
\usepackage{microtype}
\usepackage[left=0.78in,right=0.78in,top=0.76in,bottom=0.72in]{geometry}
\usepackage{setspace}
\usepackage{amsmath,amssymb,mathtools}
\usepackage{mathrsfs}
\usepackage{fancyhdr}
\usepackage{needspace}
\usepackage[hidelinks]{hyperref}
\hypersetup{
  pdftitle={How the Proof Works},
  pdfauthor={Thomas Birnie}
}

\setstretch{1.13}
\setlength{\parindent}{0pt}
\setlength{\parskip}{0.72\baselineskip}
\setlength{\abovedisplayskip}{15pt plus 3pt minus 2pt}
\setlength{\belowdisplayskip}{15pt plus 3pt minus 2pt}
\setlength{\abovedisplayshortskip}{14pt plus 2pt minus 2pt}
\setlength{\belowdisplayshortskip}{14pt plus 2pt minus 2pt}
\setlength{\jot}{9pt}
\setlength{\emergencystretch}{2em}
\allowdisplaybreaks[2]
\raggedbottom
\clubpenalty=10000
\widowpenalty=10000
\displaywidowpenalty=10000

\pagestyle{fancy}
\fancyhf{}
\renewcommand{\headrulewidth}{0pt}
\fancyfoot[C]{\small\thepage}

\newcommand{\R}{\mathbb R}
\newcommand{\I}{\mathcal I}

\begin{document}

\begin{center}
  {\LARGE\bfseries How the Proof Works}\par
  \vspace{0.35\baselineskip}
  {\normalsize Thomas Birnie}\par
\end{center}

\vspace{0.55\baselineskip}

Fix a smooth divergence-free initial velocity \(u_0\), a viscosity \(\nu>0\),
and the classical solution \((u,p)\) on its maximal interval \([0,T_*)\).
The three-dimensional incompressible Navier--Stokes equation is
\[
  \partial_tu+(u\cdot\nabla)u
  =-\nabla p+\nu\Delta u,
  \qquad
  \nabla\cdot u=0.
\]
Transport carries velocity through its own gradients, pressure keeps the field
incompressible, and viscosity diffuses those gradients. Assume for contradiction
that \(T_*<\infty\). The gradients would have to keep sharpening until the
solution lost smoothness.

That sharpening does not require infinite kinetic energy. As a vortex stretches,
incompressibility makes its core thinner. If the velocity changes by an amount
of size \(|u|\) across a core of radius \(r\), then its gradient has the scale
\[
  |\nabla u|\sim\frac{|u|}{r}.
\]
Shrinking the core raises the gradient. Yet the energy identity still gives
\[
  \frac12\|u(t)\|_2^2
  +\nu\int_0^t\|\nabla u(\tau)\|_2^2\,d\tau
  =\frac12\|u_0\|_2^2.
\]
The fluid can keep finite total energy while its motion concentrates into a
smaller region and its gradients become steeper.

To see where the Sobolev scale still detects that concentration, compare the
solution with a compressed copy of itself. For \(\lambda>1\), set
\[
  u_\lambda(x,t):=\lambda u(\lambda x,\lambda^2t),
  \qquad
  p_\lambda(x,t):=\lambda^2p(\lambda x,\lambda^2t).
\]
With \(y=\lambda x\) and \(\tau=\lambda^2t\), the four terms become
\[
\begin{aligned}
  \partial_tu_\lambda(x,t)
  &=\lambda^3(\partial_tu)(y,\tau),
  \\
  (u_\lambda\cdot\nabla)u_\lambda(x,t)
  &=\lambda^3\bigl((u\cdot\nabla)u\bigr)(y,\tau),
  \\
  \nabla p_\lambda(x,t)
  &=\lambda^3(\nabla p)(y,\tau),
  \\
  \nu\Delta u_\lambda(x,t)
  &=\lambda^3\nu(\Delta u)(y,\tau).
\end{aligned}
\]
Compression strengthens transport and diffusion by an equal factor. Scale
alone gives neither term an advantage.

The Sobolev norm separates the levels. The Fourier transform obeys
\[
  \widehat{u_\lambda}(\xi,t)
  =\lambda^{-2}\widehat u(\xi/\lambda,\tau),
\]
and the change of variables \(\eta=\xi/\lambda\) gives
\[
\begin{aligned}
  \|u_\lambda(t)\|_{\dot H^s}^2
  &=\int_{\R^3}|\xi|^{2s}\lambda^{-4}
    |\widehat u(\xi/\lambda,\tau)|^2\,d\xi
  \\
  &=\lambda^{2s-1}
    \|u(\tau)\|_{\dot H^s}^2.
\end{aligned}
\]
At \(s=0\), the squared norm shrinks by \(\lambda^{-1}\); at \(s=1\), it
grows by \(\lambda\). Exactly between them, \(2s-1=0\), so
\[
  s=\frac12,
  \qquad
  \|u_\lambda(t)\|_{\dot H^{1/2}}
  =\|u(\tau)\|_{\dot H^{1/2}}.
\]
The energy level loses size under compression, the first-derivative level gains
it, and \(\dot H^{1/2}\) stays fixed. That unchanged middle level is the critical
height.

The rescaling assigns a shorter time window to each narrower copy. Suppose one
fluid history follows that parabolic pattern and reaches a geometric sequence of
core radii:
\[
  r_n=2^{-n}r_0
\]
Parabolic scaling gives the \(n\)th stage a duration proportional to \(r_n^2\),
so halving the radius quarters the duration:
\[
  \delta t_n=4^{-n}\delta t_0,
  \qquad
  \sum_{n=0}^{\infty}\delta t_n
  =\frac43\,\delta t_0<\infty.
\]
The geometric sum permits infinitely many finer stages before a finite time
\(T_*\). Along such a history, the kinetic energy may stay finite while each
stage carries a narrower core, a steeper gradient, and a shorter duration. This
is the Zeno cascade.

\Needspace{9\baselineskip}
In order for a singularity to form, velocity has to blow up. In order for that
to happen, its acceleration would also have to blow up, which means its jerk has
to blow up, and so on, and so on:
\[
  u,
  \qquad
  \partial_tu,
  \qquad
  \partial_t^2u,
  \qquad
  \ldots
\]
A cascade of narrower cores and steeper gradients suggests less spatial
regularity as temporal order rises. Acceleration should look rougher than
velocity, jerk rougher than acceleration. Differentiating the equation reveals
the opposite.

\begingroup
\small
\begin{flalign*}
  &\partial_tu
  =\nu\Delta u-(u\cdot\nabla)u-\nabla p,&&
  \\[0.6\baselineskip]
  &\partial_t^2u
  =\nu\Delta\partial_tu-(\partial_tu\cdot\nabla)u
    -(u\cdot\nabla)\partial_tu-\nabla\partial_tp,&&
  \\[0.6\baselineskip]
  &\partial_t^3u
  =\nu\Delta\partial_t^2u-(\partial_t^2u\cdot\nabla)u
    -2(\partial_tu\cdot\nabla)\partial_tu
    -(u\cdot\nabla)\partial_t^2u-\nabla\partial_t^2p.&&
\end{flalign*}
\endgroup

The branching visible in those three lines continues at every finite order
\(n\):
\[
\begin{aligned}
  \partial_t^{\,n+1}u
  ={}&\nu\Delta\partial_t^nu
  \\
  &-\sum_{a=0}^{n}\binom{n}{a}
    \bigl(\partial_t^au\cdot\nabla\bigr)
    \partial_t^{\,n-a}u
  \\
  &-\nabla\partial_t^np.
\end{aligned}
\]
Call the \(n\)th velocity and pressure responses
\[
  V_n:=\partial_t^nu,
  \qquad
  Q_n:=\partial_t^np.
\]

Let \(\Lambda:=(-\Delta)^{1/2}\), so \(\|\Lambda^sf\|_2\) measures spatial
Sobolev height \(s\). At temporal order \(r\), the energy at that height is
\[
  E_{r,s}(t):=\frac12\|\Lambda^sV_r(t)\|_2^2
\]
The nonlinear and pressure terms contribute
\[
\begin{aligned}
  \mathcal P_{r,s}(t)
  &:=-\operatorname{Re}
  \Biggl\langle
    \Lambda^sV_r,
    \Lambda^s\!\left[
      \sum_{a=0}^{r}\binom{r}{a}
      (V_a\cdot\nabla)V_{r-a}
      +\nabla Q_r
    \right]
  \Biggr\rangle_{L^2}.
\end{aligned}
\]
Since \(\Delta=-\Lambda^2\), the viscous pairing is
\[
\begin{aligned}
  \nu\operatorname{Re}
  \langle\Lambda^sV_r,\Lambda^s\Delta V_r\rangle_{L^2}
  &=-\nu\|\Lambda^{s+1}V_r\|_2^2
  \\
  &=-2\nu E_{r,s+1}.
\end{aligned}
\]
The Laplacian places one spatial derivative on each factor, tying height \(s\)
directly to height \(s+1\):
\[
  \partial_tE_{r,s}
  =\mathcal P_{r,s}-2\nu E_{r,s+1},
  \qquad
  E_{r,s+1}
  =\frac{\mathcal P_{r,s}-\partial_tE_{r,s}}{2\nu}.
\]
One more derivative gives
\[
  \partial_t^2E_{r,s}
  =\partial_t\mathcal P_{r,s}
  -2\nu\mathcal P_{r,s+1}
  +(2\nu)^2E_{r,s+2},
\]
Iterating \(k\) times gives
\[
  \partial_t^kE_{r,s}
  =(-2\nu)^kE_{r,s+k}
  +\sum_{j=0}^{k-1}
    (-2\nu)^j
    \partial_t^{\,k-1-j}\mathcal P_{r,s+j}.
\]

\Needspace{22\baselineskip}
At critical height, set \(H:=E_{0,1/2}\). Each derivative of \(H\) is a
quadratic contraction of the velocity responses:
\[
  H^{(k)}
  =\frac12\sum_{a=0}^{k}\binom{k}{a}
  \operatorname{Re}
  \left\langle
    \Lambda^{1/2}V_a,
    \Lambda^{1/2}V_{k-a}
  \right\rangle_{L^2}.
\]
The velocity responses build the temporal derivatives of \(H\), and the viscous
identity exposes a higher spatial energy at each order:
\[
\begin{aligned}
  u,\quad \partial_tu,\quad \partial_t^2u,\quad\ldots
  &\longrightarrow
  H,\quad H',\quad H'',\quad\ldots
  \\
  &\longrightarrow
  E_{0,1/2},\quad E_{0,3/2},\quad E_{0,5/2},\quad\ldots .
\end{aligned}
\]
Each rise in temporal order exposes a rise in spatial Sobolev order. The
transfer terms remain in the formula; the top energy is isolated through the
nonzero coefficient \((-2\nu)^k\). If every finite height reaches \(T_*\),
Sobolev embedding gives, for every finite \(k\),
\[
  s>k+\frac32
  \quad\Longrightarrow\quad
  H^s(\R^3)\hookrightarrow C^k(\R^3),
\]
and hence
\[
  \bigcap_{m<\infty}H^m(\R^3)
  =H^\infty(\R^3)
  \hookrightarrow C^\infty(\R^3).
\]
The temporal sequence that was expected to lose spatial regularity has instead
opened every finite Sobolev height and, through their intersection, smoothness.

\Needspace{12\baselineskip}
To carry those heights to \(T_*\), pair each temporal response with the Sobolev
space in which it is measured. The coordinate \((n,m)\) means
\[
  V_n=\partial_t^nu
  \quad\text{measured in}\quad
  H_x^m.
\]
Choosing finite cutoffs \(N,M\) gathers the needed coordinates into one
rectangle. On a strict cutoff \(T<T_*\), its adjacent pairs are measured by
\[
\begin{aligned}
  \mathfrak R_{N,M}(u;T)
  :={}&
  \sum_{n=0}^{N}
  \|V_n\|_{L^2(0,T;H^{M+1})}^2
  \\
  &+
  \sum_{n=0}^{N}
  \|V_{n+1}\|_{L^2(0,T;H^{M-1})}^2,
\end{aligned}
\]
and the whole-terminal rectangle theorem gives
\[
  \sup_{0<T<T_*}\mathfrak R_{N,M}(u;T)<\infty.
\]
Since the integrands are nonnegative, increasing \(T\) to \(T_*\) yields
\[
  V_n\in L^2(0,T_*;H^{M+1}),
  \qquad
  V_{n+1}\in L^2(0,T_*;H^{M-1}),
  \qquad 0\leq n\leq N.
\]
With \(A:=(1-\Delta)^{1/2}\), the two rows meet at height \(M\):
\[
\begin{aligned}
  \|V_n(b)\|_{H^M}^2-\|V_n(a)\|_{H^M}^2
  =2\operatorname{Re}\int_a^b
  \bigl\langle
    A^{M-1}V_{n+1}(t),
    A^{M+1}V_n(t)
  \bigr\rangle_{L^2}\,dt.
\end{aligned}
\]
Cauchy--Schwarz gives
\[
  V_n\in C([0,T_*];H^M).
\]
The pressure recurrence supplies
\[
  Q_n\in L^1(0,T_*;H^M),
  \qquad
  0\leq n\leq N,
\]
which completes the finite velocity--pressure rectangle
\(\mathcal R_{N,M}[u_0]\).

A larger rectangle restricts to a smaller one:
\[
  \rho_{N',M'}^{N,M}\mathcal R_{N,M}[u_0]
  =\mathcal R_{N',M'}[u_0].
\]
Taking all finite rectangles together gives
\[
  \I[u_0]
  :=\bigl(\mathcal R_{N,M}[u_0]\bigr)_{N,M<\infty},
  \qquad
  u\in
  \bigcap_{r,m<\infty}
  H^r(0,T_*;H^m(\R^3)).
\]
Every finite coordinate appears in some rectangle, and the rectangles agree
wherever they overlap. The zeroth temporal row
\(\mathscr S_0:=\pi_0\I[u_0]\) contains \(u\) at every finite spatial height.

Taking \(n=0\) in the trace relation gives, for every finite \(m\),
\[
  u\in H^1(0,T_*;H^m)
  \hookrightarrow C([0,T_*];H^m).
\]
The traces converge to one endpoint, smooth at every order:
\[
  u_*:=\lim_{t\uparrow T_*}u(t)
  \in\bigcap_{m<\infty}H^m(\R^3)
  =H^\infty(\R^3)
  \hookrightarrow C^\infty(\R^3).
\]
Local smooth existence from \(u_*\) continues the solution through
\([T_*,T_*+\delta]\), and uniqueness joins it to the existing history. A finite
\(T_*\) cannot be maximal, so
\[
  T_*=\infty.
\]

Each familiar finite estimate comes from one coordinate of the intersection.
Given a finite temporal order \(r\), spatial derivative order
\(k\), and \(2\leq p\leq\infty\), choose \(m\) so that
\[
  H^m(\R^3)\hookrightarrow W^{k,p}(\R^3).
\]
When \(p=\infty\), any \(m>k+\tfrac32\) is enough; for finite \(p\), choose
\(m\) in the standard Sobolev range. That coordinate gives
\[
  \partial_t^ru\in C_tH_x^m
  \quad\Longrightarrow\quad
  \nabla^k\partial_t^ru\in L_t^\infty L_x^p.
\]
For example,
\[
\begin{aligned}
  \I[u_0]&\longrightarrow C_tH_x^1
    \longrightarrow L_t^\infty L_x^6,
  \\
  \I[u_0]&\longrightarrow C_tH_x^2
    \longrightarrow L_t^\infty L_x^\infty,
  \\
  \I[u_0]&\longrightarrow C_tH_x^3
    \longrightarrow L_t^\infty W_x^{1,\infty}.
\end{aligned}
\]
The last line, for instance, gives on every finite interval \([0,T]\)
\[
\begin{aligned}
  \int_0^T\|\nabla u(t)\|_\infty\,dt
  &\leq T\sup_{0\leq t\leq T}\|\nabla u(t)\|_\infty
  \\
  &\leq C T\sup_{0\leq t\leq T}\|u(t)\|_{H^3}
  <\infty.
\end{aligned}
\]
The \(H^3\) line is one instance. Changing \(r\), \(k\), or \(p\) only changes
which finite height is taken from the intersection. The order is intersection
first, estimate second.

In shorthand:
\[
\begin{aligned}
  u_0
  &\longrightarrow \{(V_n,Q_n)\}_{n\geq0}
  \longrightarrow \{\mathcal R_{N,M}[u_0]\}_{N,M<\infty}
  \\
  &\longrightarrow \I[u_0]
  \longrightarrow \mathscr S_0
  \longrightarrow u_*\in H^\infty
  \longrightarrow T_*=\infty.
\end{aligned}
\]

\end{document}
